% Section 2: Related Work
\section{Related Work}
\label{sec:related}

\subsection{Verbosity Bias and Length-Controlled Evaluation}

The problem of length bias in LLM-based evaluation is well-documented.
\citet{Zheng2023Judging} show that GPT-4 as a pairwise judge prefers longer, more structured
responses even when content quality is equal. \citet{Shi2024Judging} extend this to systematic
position and length bias across judge models.

\citet{Dubois2024LengthControlled} respond with AlpacaEval 2.0's length-controlled win rate
(\texttt{LC\_winrate}). Their GLM-based approach regresses response length out of pairwise
comparison outcomes, yielding a debiased preference score. They report a bivariate Spearman
correlation of $\rho \approx 0.94$ between \texttt{win\_rate} and \texttt{LC\_winrate}, and
show \texttt{LC\_winrate} correlates more strongly with Chatbot Arena ($\rho \approx 0.98$).
Our work extends this: we move from bivariate to partial correlation, controlling for verbosity
(\texttt{avg\_length}), and find the capability-LC alignment strengthens to 0.985 when
verbosity is properly held constant.

\citet{Hu2024Explaining} provide a mechanistic decomposition of \texttt{win\_rate} into a
desirability component (length-independent quality) and information mass component
(length-dependent content). Their framework predicts the desirability channel survives LC
correction, consistent with our finding that $\beta_{\mathrm{win}}$ dominates
$\beta_{\mathrm{len}}$ in standardized regression.

\subsection{Bidirectional Human-AI Alignment}

\citet{Shen2024Bidirectional} provide a systematic review of bidirectional alignment between
humans and AI systems, identifying a gap in empirical quantification of capability-modulated
alignment asymmetry. \citet{Ji2023Survey} survey alignment approaches broadly, noting that
evaluation metric consistency across human and AI evaluators is an open problem. Our work
directly addresses the empirical gap identified by \citet{Shen2024Bidirectional}: we quantify,
at population scale, how strongly capability modulates the alignment between human preference
(\texttt{win\_rate}) and AI-debiased preference (\texttt{LC\_winrate}).

\subsection{Positioning}

Our work differs from prior studies in three ways: (1) we study the \textit{partial}
relationship between capability and LC preference, explicitly controlling for verbosity;
(2) we provide a FWL-based methodological robustness check new to this literature; and
(3) we analyze both \texttt{LC\_winrate} and $\Delta$ as dependent variables, providing a
comparative effect-size lesson with implications for evaluation study design. We complement the
AlpacaEval 2.0 methodology: our results validate the LC correction as a
capability-preserving transformation.
